% Resumen de la aplicación web (Front-end)

\section{Descripción general del front-end}

El front-end de la aplicación está implementado con el framework \textbf{Angular} (versión 19) y utiliza la plantilla de panel de administración \textbf{CoreUI} sobre \textit{Bootstrap}. Este front-end proporciona una interfaz web responsiva orientada a aplicaciones tipo tablero de control (\emph{dashboard}), que sirve como punto de acceso para la visualización y gestión de la información generada por el sistema.

La estructura general se basa en un \emph{layout} principal, denominado \texttt{DefaultLayoutComponent}, que define la barra lateral de navegación, la cabecera y el área de contenido central. Todas las vistas funcionales de la aplicación se cargan de forma perezosa (\emph{lazy loading}) como módulos hijos de este layout, lo que mejora el rendimiento y la escalabilidad del proyecto.

\section{Arquitectura y organización de módulos}

El enrutamiento principal se define en el archivo \texttt{app.routes.ts}. Desde este punto se configuran las rutas que conforman la interfaz de usuario:

\begin{itemize}
  \item \textbf{Ruta raíz} (``/``): redirige automáticamente al módulo de tablero de control en ``/dashboard``.
  \item \textbf{Módulo Dashboard} (``/dashboard``): vista principal del sistema, pensada para mostrar indicadores, gráficas y widgets de estado general.
  \item \textbf{Módulo de componentes base} (``/base``): agrupa componentes de interfaz como tarjetas, tablas, acordeones, paginación, etc., que sirven de soporte para la presentación de datos.
  \item \textbf{Módulo de botones} (``/buttons``): recopila distintos estilos de botones y grupos de botones utilizados en la navegación y en los formularios.
  \item \textbf{Módulo de formularios} (``/forms``): contiene los formularios de entrada de datos (campos de texto, selectores, validaciones, etc.), que permiten al usuario interactuar con la información del sistema.
  \item \textbf{Módulo de iconos} (``/icons``): proporciona la galería de iconos utilizados en la interfaz para mejorar la usabilidad y la identificación visual de acciones.
  \item \textbf{Módulo de notificaciones} (``/notifications``): muestra ejemplos de alertas, mensajes y toasts para retroalimentación al usuario.
  \item \textbf{Módulo de widgets} (``/widgets``): incluye bloques de información reutilizables que pueden integrarse en el tablero u otras vistas.
  \item \textbf{Módulo de gráficas} (``/charts``): implementa visualizaciones con \texttt{chart.js} integradas con CoreUI (gráficas de líneas, barras, etc.) para la representación visual de datos.
  \item \textbf{Módulo de páginas} (``/pages``): agrupa páginas especiales como error 404, error 500, pantalla de inicio de sesión y registro de usuarios.
\end{itemize}

Adicionalmente, se definen rutas explícitas para las páginas de error (``/404`` y ``/500``), así como para las pantallas de autenticación: ``/login`` y ``/register``. Esto permite gestionar de forma clara los flujos de acceso al sistema y el manejo de errores de navegación.

\section{Navegación y experiencia de usuario}

La navegación lateral se configura en el archivo \texttt{_nav.ts}, donde se declaran los elementos del menú (\texttt{navItems}). Cada elemento se asocia con una ruta y un icono, integrando la biblioteca de iconos de CoreUI. El menú incluye al menos la entrada \textit{Dashboard}, y puede ampliarse fácilmente con nuevas secciones vinculadas a funcionalidades específicas del dominio de la aplicación.

El diseño es responsivo y está construido sobre componentes reutilizables. Esto facilita mantener una apariencia consistente entre las distintas vistas y reduce el esfuerzo de desarrollo al agregar nuevas pantallas.

\section{Servicios y visualización de datos}

El front-end incorpora servicios Angular, como \texttt{dashboard-data.service.ts}, que encapsulan la obtención y preparación de datos para las vistas del tablero. Estos servicios pueden conectarse a una API REST en el back-end de la aplicación para recuperar información en tiempo real (por ejemplo, estadísticas, series temporales, resultados de modelos, etc.).

Para la visualización gráfica se emplea la librería \texttt{chart.js} junto con los componentes de \texttt{@coreui/angular-chartjs}. Esto permite representar indicadores mediante gráficos de líneas, barras, áreas u otros tipos, lo cual es adecuado para mostrar tendencias y comparaciones en el contexto de la tesis.

\section{Interacción con el back-end}

Aunque el repositorio del front-end está desacoplado del back-end, la arquitectura está preparada para consumir servicios REST mediante \texttt{HttpClient}. De este modo, el front actúa como capa de presentación sobre la lógica y los datos expuestos por el servidor (por ejemplo, un API en Python/FastAPI). Esta separación facilita la evolución independiente de cada parte y permite reutilizar la misma API para otros clientes (aplicaciones móviles, scripts de análisis, etc.).

\section{Rol del front-end en la tesis}

En el contexto de la tesis, este front-end cumple el papel de \textbf{interfaz de usuario} del sistema propuesto. Sus principales aportes son:

\begin{itemize}
  \item Proporcionar una plataforma unificada para la consulta de resultados, indicadores y estadísticas generados por el modelo o sistema desarrollado en el back-end.
  \item Ofrecer mecanismos de interacción (formularios, filtros, navegación entre módulos) que permiten al usuario explorar los datos de forma intuitiva.
  \item Presentar los resultados mediante visualizaciones gráficas y tablas, facilitando la interpretación y el análisis por parte de usuarios no expertos.
  \item Servir como demostración práctica de la aplicabilidad del sistema, al integrar los componentes técnicos en una herramienta web accesible.
\end{itemize}

Este resumen puede incluirse en el capítulo de diseño o implementación de la tesis para describir la capa de presentación de la solución, detallando tanto las tecnologías empleadas como la organización funcional de la interfaz web.
